\section{Soundness}

\begin{lemma}[Good challenges keep the agreement]\label{lem:sound-fold}
Assume $\mathrm H(\alpha_*,\Theta)$ with $\alpha_*>\rho$, let $\alpha\ge\alpha_*$, and let
$\agree(u_0,C_{w,D})\le\alpha$. Fix a prover and challenges $r_1,\dots,r_m\in\F$, let
$u_1,\dots,u_m$ be the prover's words, and let $S_i\subseteq\F$ be the set of \cref{lem:round}
for $u:=u_{i-1}$ and $\mu^{(i-1)}$; it depends only on $r_1,\dots,r_{i-1}$. If $r_i\notin S_i$
for every $1\le i\le m$, then $\agree_{\mu^{(m)}}(u_m,\Gc_m^{w,D})\le\alpha$.
\end{lemma}

\begin{proof}
By \cref{lem:start}, $\agree_{\mu^{(0)}}(u_0,\Gc_0^{w,D})=\agree(u_0,C_{w,D})\le\alpha$. For
$1\le i\le m$, $r_i\notin S_i$ means by \cref{lem:round} that
$\agree_{\nu^{(i)}}(\Fold_{c_i}(u_{i-1}),\Gc_i^{w,D})\le\beta_i$, and $\beta_i$ is the maximum of
$\agree_{\mu^{(i-1)}}(u_{i-1},\Gc_{i-1}^{w,D})$ and $\alpha_*$; with \cref{lem:no-increase} this
gives
$\agree_{\mu^{(i)}}(u_i,\Gc_i^{w,D})\le\max(\agree_{\mu^{(i-1)}}(u_{i-1},\Gc_{i-1}^{w,D}),\alpha_*)$,
and induction from the base case, using $\alpha\ge\alpha_*$, gives the claim.
\end{proof}
\uses{\cref{hyp:ca,lem:start,lem:round,lem:no-increase}}\usedby{\cref{thm:reduction}}

\begin{theorem}[Soundness, reduction form]\label{thm:reduction}
Assume $\mathrm H(\alpha_*,\Theta)$ with $\alpha_*>\rho$, let $\alpha\ge\alpha_*$, and suppose
$\agree(u_0,C_{w,D})\le\alpha$. Then for every prover the test $\Pi_{w,D}$ accepts with
probability at most
\[
\frac1{|\F|}\sum_{i=1}^{m}\Theta\big(|L_i|,k^i,k-1\big)+\alpha^s .
\]
\end{theorem}

\begin{proof}
The prover's words are fixed; the challenges $(r_1,\dots,r_m)$ are uniform on $\F^m$. By
\cref{lem:bad-count} with $S_i$ as in \cref{lem:round} and $\Theta_i:=\Theta(|L_i|,k^i,k-1)$, the
tuples with $r_i\in S_i$ for some $i$ number at most $|\F|^{m-1}\sum_i\Theta_i$, so they have
probability at most $\frac1{|\F|}\sum_i\Theta_i$. Outside them \cref{lem:sound-fold} gives
$\agree_{\mu^{(m)}}(u_m,\Gc_m^{w,D})\le\alpha$; since $u_m=g\in\Gc_m^{w,D}$ and $\Gc_m^{w,D}$ is the
set of constant words with value in $\F[T]_{\le D}$, \cref{ext:claim85} bounds the probability that
one query accepts by $\alpha$ and the $s$ repetitions by $\alpha^s$. Adding the two bounds gives the
claim.
\end{proof}
\uses{\cref{lem:sound-fold,lem:bad-count,ext:claim85}}\usedby{\cref{cor:johnson} and \cref{cor:ud}}

\begin{consequence}[Soundness at the Johnson radius]\label{cor:johnson}
Let $\kappa\ge3$, $\alpha\ge\alpha_0:=\sqrt\rho\,(1+\frac1{2\kappa})$, and
suppose $\agree(u_0,C_{w,D})\le\alpha$. Then the test accepts with probability at most
$\varepsilon_C+\alpha^s$, where
\[
\varepsilon_C:=\frac1{|\F|}\Big(\frac{(1+\frac1{2\kappa})^7\kappa^7\,N^2}{3(k+1)\rho^{3/2}}
+m(k-1)\frac{(2\kappa+1)(N+1)}{\sqrt\rho}\Big).
\]
\end{consequence}

\begin{proof}
Use \cref{prop:H1} in \cref{thm:reduction}: for each $i$ the quantity
$\Theta_J(|L_i|,k^i,k-1)$ equals $(k-1)\max(a_J|L_i|^2,b_J(N+1))$, with $a_J$ and $b_J$ the two
constants of $\Theta_J$. Here $M=k^i$, so $Mn=N$. Hence
\[
\sum_{i=1}^{m}\Theta_J(|L_i|,k^i,k-1)\ \le\ a_J\sum_{i=1}^{m}(k-1)|L_i|^2+b_J\,m(k-1)(N+1),
\]
and $\sum_i(k-1)|L_i|^2=(k-1)N^2\sum_ik^{-2i}=\frac{N^2(1-k^{-2m})}{k+1}\le\frac{N^2}{k+1}$.
\end{proof}
\uses{\cref{prop:H1,thm:reduction}}\usedby{the paper}

\begin{consequence}[Soundness in the unique-decoding regime]\label{cor:ud}
Let $\alpha>\frac{1+\rho}2$ and suppose $\agree(u_0,C_{w,D})\le\alpha$. Then the test accepts
with probability at most $\varepsilon_{\mathrm U}+\alpha^s$, where
\[
\varepsilon_{\mathrm U}:=\frac{m(k-1)(N+1)}{|\F|}\Big(\frac1{\eta_{\mathrm U}}+\frac3{\sqrt\rho}\Big),
\qquad \eta_{\mathrm U}=\min\Big(\frac{(1-\sqrt\rho)^2}2,\frac{\sqrt\rho}{20}\Big).
\]
\end{consequence}

\begin{proof}
Use \cref{prop:H2} with $\alpha_*=\alpha$ in \cref{thm:reduction}: each of the $m$ terms
$\Theta_U(|L_i|,k^i,k-1)$ equals $(k-1)(N+1)(\frac1{\eta_{\mathrm U}}+\frac3{\sqrt\rho})$ by \cref{def:domains}.
\end{proof}
\uses{\cref{prop:H2,thm:reduction}}\usedby{the paper}

\begin{remark}[Individual degrees]
Since $C_{w,D}\subseteq\RS[\F,L_0,k^m]$, acceptance also certifies $\deg_{x_l}<k$ for all $l$;
taking $w=e_l$ and $D=b<k$ certifies $\deg_{x_l}f\le b$.
\end{remark}

\begin{remark}[Characteristic]\label{rem:char}
\Cref{prop:H1} inherits from \cref{ext:bciks72} the separability convention of
[BCIKS20, \S5.2.3] (the characteristic of $\F$ exceeds $D_Y$, the bound on the $Y$-degree of the interpolating polynomial $Q(X,Y,Z)$ of [BCIKS20, \S5], which satisfies $D_Y<\frac{2\kappa+1}{2\sqrt\rho}$); the
general case, with inseparable factors, is treated in [BCIKS20, Appendix~C].
\end{remark}
\uses{\cref{ext:bciks72}}\usedby{the paper}
